% Reproducibility details; the main argument does not depend on reading this appendix.
\chapter{Software interfaces and reproduction notes}
\label{app:technical-reference}

This appendix records the concrete software interfaces and historical examples
behind the main narrative. It is intended for someone reproducing or extending
the implementation. The mathematical definitions, worked decisions and proof
assumptions remain in the main chapters.

% BEGIN SOURCE UNIT P-03d-java-delivery:00
\section{The earlier SPI compiler: a concrete deployment lesson}

The API and commands in the next five subsections belong to the retained
\emph{SPI-Demo1} teaching example under \path{examples/java-dry-run/}. They remain
reproducible, but are not the current v0.1 SDK. The subsequent section on the actual
generated interface supplies the current String/Map API. Keeping both makes the
translation mechanism inspectable without directing an implementer to the wrong
library.


\label{merge:P-03d-java-delivery:00}

\label{sec:java}

Java is the required deployment output. The authoring workbench may use Python
for source processing, review and reference evaluation; a transaction must be
able to call the released Java library without running Python or contacting a
language model. The initial integration assumption is a plain Java 17 JAR.
This keeps the decision code independent of Spring, application servers and
particular rule engines. A bank on another Java version needs a declared target
and the same conformance checks before receiving a compatible build.
% END SOURCE UNIT P-03d-java-delivery:00

% BEGIN SOURCE UNIT P-03d-java-delivery:01
\subsection{SPI-Demo1: the narrower retained demonstration}
\label{merge:P-03d-java-delivery:01}


The delivered \texttt{SPI-Demo1} example contains a source-anchored RuleIR bundle,
synthetic fact snapshots, a separately written Python evaluator, a deterministic
Java emitter, the generated class, two small Java support classes, and a separate
host caller. Its implemented expression subset is deliberately explicit:
Boolean, integer and HK-cent literals and facts; unconditional rule references;
\texttt{all}, \texttt{any}, \texttt{not}; and exact numeric
\texttt{eq}/\texttt{ge}/\texttt{gt} comparisons. The root has a Boolean scope.
The demo compiler rejects dates, arithmetic expressions, general defaults,
conditionals and the full obligation language. The later RuleIR implementation
supports the specified decision operators and its finite duty profile; those
features are outside this earlier demo API. The demonstration response is its own smaller API profile,
not a claim to implement every field and trace node of the full RuleIR service.

% BEGIN REVISION ADDITION teach-P-03d-java-delivery-01
\begin{figure}[H]
\centering
\TeachingCompare{SPI-Demo1's limited operators}{Full RuleIR runtime}{Different API and evidence sets}
\caption{The early compiler example and the later runtime have distinct scopes.}\label{fig:teach-P-03d-java-delivery-01}
\end{figure}

% END REVISION ADDITION teach-P-03d-java-delivery-01
The distinction matters for an implementing agent. The example demonstrates an
actual compiler and integration path that can be extended. It does not silently
turn the 35 general RuleIR contract fixtures into executed tests. Those fixtures were the completion target for the full interpreter and
backend and were subsequently executed by the accepted MVP. The separate
SPI example has 32 executed decision scenarios and eleven consent histories, with
its own recorded evidence.
% END SOURCE UNIT P-03d-java-delivery:01

% BEGIN SOURCE UNIT P-03d-java-delivery:02
\subsection{The SPI-Demo1 input and output boundary}
\label{merge:P-03d-java-delivery:02}


The public API accepts an immutable \texttt{Snapshot}, a time being assessed and
a knowledge cutoff. A \texttt{Snapshot} binds a subject identifier to the
declared fact map. A \texttt{Fact} carries a declared scalar type and one of
\texttt{known}, \texttt{unknown} or \texttt{conflict}. A known Boolean is a Java
\texttt{Boolean}; a known monetary or count value is a \texttt{BigInteger}.
Its evidence identifiers, validity interval and recording time are mandatory.
Unknown has a reason and no value. Conflict identifies competing evidence.

The host supplies every declared field, using an explicit unknown where evidence
is unavailable. Omitting a field or supplying the wrong type is a malformed
request, not a legal conclusion. Validity intervals are half-open: evidence
valid until 10:00 does not establish a fact at 10:00. Evidence recorded after
the knowledge cutoff is unavailable to that replay. The API takes canonical UTC
timestamps with six fractional digits; conversion from a Hong Kong business
timestamp belongs in the host adapter and must preserve the intended instant.

\begin{lstlisting}[language=Java,caption={The library API and one dated monetary input.}]
DecisionRuntime.Result evaluate(
    DecisionRuntime.Snapshot snapshot,
    String validAt,
    String knownAt);

// A known value is exact and supported by dated evidence.
Fact portfolio = Fact.known(
    "money_hkd", new BigInteger("4000000000"),
    List.of("valuation.lee", "ownership.lee"),
    "2026-09-01T00:00:00.000000Z",
    "2026-10-01T00:00:00.000000Z",
    "2026-09-01T00:00:00.000000Z");

Fact missingPortfolio = Fact.unknown("money_hkd", "MISSING");
\end{lstlisting}
The first amount is HK\$40 million expressed in cents. A bank adapter reading a
decimal amount should specify the currency, use exact decimal arithmetic and
reject an unsupported fraction of a cent after the adopted conversion. In Java,
\texttt{amount.movePointRight(2).toBigIntegerExact()} can enforce exact cents
after that policy has been applied. Calling \texttt{doubleValue()} first would
discard the property the specification is trying to preserve. For a currency
conversion, the exchange-rate source, date, direction and rounding rule are part
of the evidence mapping.

% BEGIN REVISION ADDITION teach-P-03d-java-delivery-02
\begin{figure}[H]
\centering
\TeachingFlow{Exact dated Snapshot}{SPI-Demo1 evaluation}{Tagged result and disposition}
\caption{The earlier Java interface illustrates the input and output boundary.}\label{fig:teach-P-03d-java-delivery-02}
\end{figure}

% END REVISION ADDITION teach-P-03d-java-delivery-02
The result contains a logical status and an operational disposition:
\begin{center}\small
\begin{tabularx}{\textwidth}{P{0.18\textwidth}P{0.34\textwidth}Y}
\toprule Status & Disposition & Host meaning \\\midrule
\texttt{TRUE} & Conditions met & Selected streamlining conditions established; continue the other controls.\\
\texttt{FALSE} & Standard process required & At least one required condition is refuted; do not apply this streamlined route.\\
\texttt{UNKNOWN} & Review required & Evidence can change the conclusion; show blocking inputs.\\
\texttt{CONFLICT} & Review required & Resolve inconsistent evidence; preserve both records.\\
\texttt{OUT\_OF\_SCOPE} & Outside profile & Select another implemented profile or the standard process.\\
\bottomrule
\end{tabularx}
\end{center}
The demo's JSON includes \texttt{profile}, \texttt{authority}, both timestamps,
the bundle and snapshot hashes, a postorder rule trace, missing and blocking
input lists, and the result hash. A rule-trace entry gives its identifier,
status, source-anchor identifiers and blocking inputs. Source anchors resolve
through the bundle included with the release. Known values and their evidence
remain in the preserved snapshot. The current RuleIR implementation provides expression-node traces,
standardized diagnostics and engine identity through the different v0.1 API
documented later in this chapter.

A malformed request throws an explicit Java input exception in this small
library. The production adapter must map that exception, a timeout or an
unavailable release to a technical-failure result and a review/standard-process
path. It must never catch an exception and substitute a successful Boolean.
The demonstration result always carries \texttt{DEMONSTRATION\_ONLY}; production
authority requires a different, approved release record and enforcement by the
host deployment process.
% END SOURCE UNIT P-03d-java-delivery:02

% BEGIN SOURCE UNIT P-03d-java-delivery:03
\subsection{How the compiler emits final Java rather than a suggestion}
\label{merge:P-03d-java-delivery:03}


The emitter first validates the JSON schema, references, identifiers, types and
acyclic rule graph. It then assigns an internal Java method name to each rule
and recursively emits its expression. Identifiers are validated data, never
pieces of executable text supplied by a language model. The generated monetary
comparison is represented by actual Java calls such as:
\begin{lstlisting}[language=Java,caption={The generated financial method's essential body. The complete class is built and tested in the supplied example.}]
return c.rule("spi.financial", sourceIds, () ->
    any(
        compare("ge", c.fact("portfolio"),
            Value.of(new BigInteger("4000000000"))),
        compare("ge", c.fact("net_assets_ex_home"),
            Value.of(new BigInteger("8000000000")))
    ));
\end{lstlisting}
\texttt{compare} returns a known Boolean when both operands are known, or an
unknown with its blocking input names. \texttt{any} implements the truth table
in Equation~\eqref{eq:or}. Java evaluates its arguments before calling the
method, so both branches have the specified eager evaluation and evidence
collection. Replacing this with Java's short-circuit \texttt{||} would change
the trace and missing-input behavior even where the final Boolean agreed.
The context memoizes referenced rules and records their results in a stable
order. A precheck detects conflicting facts in this profile's dependency set
before normal expression evaluation.

% BEGIN REVISION ADDITION teach-P-03d-java-delivery-03
\begin{figure}[H]
\centering
\TeachingFlow{Typed comparison in RuleIR}{BigInteger comparison in Java}{Same inclusive integer ordering}
\caption{The monetary comparison preserves meaning only if the input mapping preserves units.}\label{fig:teach-P-03d-java-delivery-03}
\end{figure}

% END REVISION ADDITION teach-P-03d-java-delivery-03
The production emitter should implement one mapping per supported operator,
with an explicit error for every unsupported construct. It must not ask an LLM
to write new Java on each regulatory change. The LLM can propose the typed
specification; after review, a deterministic compiler performs the mechanical
translation. This concentrates semantic review on a small language and makes
the generation step reproducible.

For the money comparison, the preservation argument is short. The IR represents
both amounts as mathematical integers of HK cents. The Java boundary constructs
the same integers as \texttt{BigInteger} values. The emitted comparison returns
true exactly when \texttt{compareTo} is nonnegative, which expresses the same
inclusive ordering. The argument depends on the input mapping preserving those
integers. It would fail if the adapter truncated decimals, mixed currencies or
overflowed a narrower integer before constructing the value. For a complete
compiler, analogous local arguments must be composed for every operator,
unknown/error case, rule reference and trace requirement. The delivered tests
check examples of preservation; they are not that universal proof.
% END SOURCE UNIT P-03d-java-delivery:03

% BEGIN SOURCE UNIT P-03d-java-delivery:04
\subsection{Reproduce the build and see the host result}
\label{merge:P-03d-java-delivery:04}


The repository retains a project-local Temurin JDK 17 for this check, verified
against the official archive checksum. It does not change the machine's default
Java installation. From the repository root, the single build command is:

\begin{minipage}{\textwidth}
\begin{lstlisting}
python3 examples/java-dry-run/build_and_verify.py \
  --jdk .localresources/java-toolchain/jdk-17.0.20.1+1
\end{lstlisting}
\end{minipage}

The Python authoring check uses the already available JSON Schema package. The
generated Java library itself has no third-party runtime dependencies. The build
uses \texttt{javac --release 17 -Xlint:all -Werror}; creates the JAR; compiles
the separate fixture and host callers against that JAR; compares complete
reference and Java results; and runs the three mutations. The package has fixed
ZIP timestamps to make repeated builds with the same toolchain reproducible.
The verification record stores compiler identity, source/compiler/JAR digests,
the executed command, results and limitations.

% BEGIN REVISION ADDITION teach-P-03d-java-delivery-04
\begin{figure}[H]
\centering
\TeachingFlow{Compile source and runtime}{Verify the exact JAR}{Call it from a separate host}
\caption{The demonstration crosses the actual Java library boundary.}\label{fig:teach-P-03d-java-delivery-04}
\end{figure}

% END REVISION ADDITION teach-P-03d-java-delivery-04
The delivered files below are operational copies of the design, not additional
reading needed to understand why it works:
\begin{center}\small
\begin{tabularx}{\textwidth}{P{0.45\textwidth}Y}
\toprule File under \path{examples/java-dry-run/} & Purpose \\\midrule
\path{spec/spi-control.bundle.json} & Source-linked typed interpretation for the chosen profile.\\
\path{spec/decision-cases.json}; \path{spec/consent-cases.json} & Independent expected statuses and synthetic evidence/history inputs.\\
\path{reference.py}; \path{compile_java.py} & Separate reference evaluation and deterministic Java generation.\\
\path{generated/GeneratedSpi.java} & Final generated decision class.\\
\path{src/main/java/hk/legalmath/spi/} & Evidence, result, three-valued arithmetic/logic support and consent replay.\\
\path{build/spi-controls-demo.jar} & Compiled library callable by a Java host.\\
\path{generated/BankHostExample.java} & Complete synthetic host caller using the public API.\\
\path{build/verification.json}; \path{build/host-output.txt} & Executed validation and host results.\\
\bottomrule
\end{tabularx}
\end{center}

The separate host prints these dispositions for the synthetic history:
\begin{lstlisting}
BEFORE_WITHDRAWAL STREAMLINING_CONDITIONS_MET
AFTER_WITHDRAWAL STANDARD_PROCESS_REQUIRED
RESULT_AUTHORITY DEMONSTRATION_ONLY
\end{lstlisting}
This is the concrete answer to how final code reaches Java: the bank depends on
the released library, constructs the declared evidence snapshot through its
adapter, calls the generated class and handles its structured result. The
compiler, public circulars and drafting tools are build-time/review-time
components. They need not be deployed into the bank's order-processing process.
% END SOURCE UNIT P-03d-java-delivery:04

% BEGIN SOURCE UNIT M07:09
\section{The actual generated Java interface}
\label{merge:M07:09}


The current emitter validates the bundle, canonicalizes it and embeds its exact
bytes in a generated class. The class name is
\texttt{hk.legalmath.Policy\_} followed by the first twenty hexadecimal digits of
the bundle hash. The class exposes the full bundle hash and delegates evaluation
to the Java semantic runtime. This is a generated, immutable policy binding; it
is not the older proposal's imagined record-based API or a proof that each legal
rule has been translated to hand-optimized Java control flow.

The two current entry points are:
\begin{lstlisting}[language=Java,caption={Implemented generated-policy API.}]
public static String evaluate(
    String snapshotJson,
    String ruleId,
    String validAt,
    String knownAt,
    String mode);

public static Map<String,Object> evaluate(
    Map<String,Object> snapshot,
    String ruleId,
    String validAt,
    String knownAt,
    String mode);
\end{lstlisting}

The string overload returns strict JSON. The map overload returns a deeply
immutable result. The generated policy and Java runtime operate without Python,
a language model, a solver or network retrieval. This property is important for
a bank integration: the interpretation workbench prepares and verifies a package,
while production decision execution can remain local and deterministic.

A concrete host compiles against the actual generated class name from the build
manifest. It supplies a validated snapshot, exact rule identifier, both assessment
times and the requested mode. The mode does not grant authority. A caller cannot
turn a draft into an approved control by passing the word ``production.'' The
release service and host must establish approval and applicability separately.

% BEGIN REVISION ADDITION teach-M07-09
\begin{figure}[H]
\centering
\TeachingFlow{Generated class binds bundle bytes}{String or Map entry point}{Immutable semantic result}
\caption{The current generated policy delegates to the verified runtime.}\label{fig:teach-M07-09}
\end{figure}

% END REVISION ADDITION teach-M07-09
The conceptual integration looks as follows. The symbolic class name in the
listing is replaced by the actual generated name when the integration fixture is
built; the release lookup and transaction methods are host responsibilities.
\begin{lstlisting}[language=Java,caption={Host responsibilities around deterministic evaluation.}]
ApprovedRelease release = releases.resolve(profile, validAt, knownAt);
release.verifyAuthorityAndDigests();
ConsistentSnapshot s = adapter.readSnapshotWithVersions(subject, order);
String resultJson = Policy_HASH.evaluate(
    s.strictJson(), release.ruleId(), validAt, knownAt, "production");
Decision decision = decoder.decodeStrictly(resultJson);
HostDisposition disposition = hostPolicy.map(decision);
transactions.commitIfVersionsUnchanged(
    s.versions(), order.id(), order.payloadHash(),
    release.identity(), decision, disposition);
\end{lstlisting}

\texttt{ApprovedRelease}, \texttt{ConsistentSnapshot} and the host transaction
interfaces in this illustration are proposed integration concepts, not classes
already supplied by the generated SDK. Their responsibilities must be implemented
in the bank's environment. The example makes the boundary explicit so an agent
does not mistake illustrative types for an existing dependency.

The host must handle every tagged result. A true selected prohibition condition
may require stopping the relevant offer; false permits only progression to other
controls under the host policy. Unknown, conflict, out-of-scope and error each
have explicit dispositions. No result string should be cast to a Boolean by a
language's truthiness rules. The decoder rejects unknown status tags so that a
future runtime extension cannot silently take an old default branch.
% END SOURCE UNIT M07:09

% BEGIN SOURCE UNIT M05:10
\section{What can be reused from the adjacent projects}
\label{merge:M05:10}


DynareMCP contains useful patterns for explicit hypotheses, diagnostics,
falsifiers and limits of claims. In the inspected legacy audit writer,
\srcpath{write_hypothesis_node} records a hypothesis, source support, code support,
a diagnostic, a falsifier and a supplied rank score. These fields map naturally
to interpretation candidates. The supplied score is not automatically calibrated;
its meaning depends on the caller.

The inspected legacy scheduler chooses among eligible frontier nodes by a stored
priority and identifier. It is not a general UCT engine hidden behind the word
``tree.'' Reusing its bookkeeping ideas is sensible; importing it and claiming
that a proven Monte Carlo search now verifies legal meaning would be wrong. The
new controller needs its own issue identity, source lineage, legal review states
and resource accounting.

% BEGIN REVISION ADDITION teach-M05-10
\begin{figure}[H]
\centering
\TeachingFlow{Adjacent project's function}{Explicit local assumptions}{Bounded adapter check}
\caption{Reuse must name a concrete capability and the changed context.}\label{fig:teach-M05-10}
\end{figure}

% END REVISION ADDITION teach-M05-10
MathDevMCP is useful once a proposition is precise enough to prove. A formal
obligation might concern preservation of a Boolean expression, absence of an
unbounded loop in a finite controller model, or equivalence of two encodings over
a declared domain. ResearchAssistant supports retrieval, parsing and retention
of the literature and regulatory evidence. Neither project supplies a privileged
interpretation of English. They strengthen different links in the chain.

The design that follows therefore combines several structures: a source inventory
to detect missing material, a candidate set to expose choices, an argument graph
to record reasons, a bounded search to investigate questions, and formal checks
to compare executable consequences. The diversity is useful because each
structure can reveal mistakes that the others overlook. Its value will depend on
how discrepancies move between them, which is the central subject of the next
chapter.
% END SOURCE UNIT M05:10

% BEGIN SOURCE UNIT P-04-prototype:03
\subsection{Adapter boundaries and the earlier reuse assessment}

The earlier repository inspection contributes the adapter requirements below.
The preceding analysis identifies the actual search-controller differences;
references here to historical UCB-style investigations do not override that
finding or establish that a generic MCTS service is already available.


\label{merge:P-04-prototype:03}


The adjacent projects supply reusable patterns and potential adapters. They are
not interchangeable components of a legal reasoner. Their local code and
documentation were inspected at the revisions recorded in the source manifest
\citep{localrepos}. The proposal makes no change to those repositories and does
not assume that their current interfaces implement the proposed legal semantics.

ResearchAssistant is useful immediately for source intake, local storage, parsing,
discovery and structured review records. It was used in this investigation to
discover literature and parse retained PDFs. The available extraction route used
\texttt{pdftotext} and reported low confidence requiring manual review. LegalMath
needs a regulatory layer above it: circular identifiers, annex relationships,
paragraph anchors, versions, replacement links and an applicability inventory.
PDF extraction quality must be assessed on the actual source layout, especially
footnotes and tables.

MathDevMCP is useful after a property has been formalized. Its inspected Lean
route binds checks to an exact target and assumptions, and its result handling
distinguishes diagnostic evidence from proof. Its source-bound obligations and
resumable evidence structures can support the question in
Equation~\eqref{eq:refinement}. The legal adapter must define the rule language's
semantics, numeric and date operations, and the exact theorem requested. Structural
agreement between a formula and a code fragment is not a substitute for semantic
equivalence. The inspected project describes an internal research status and
has its own licensing terms; reuse within a bank needs the appropriate software
and ownership review.

% BEGIN REVISION ADDITION teach-P-04-prototype-03
\begin{figure}[H]
\centering
\TeachingCompare{ResearchAssistant: source processing}{MathDevMCP: scoped mathematics}{DynareMCP: bounded investigations}
\caption{Adjacent tools contribute different capabilities, each requiring local validation.}\label{fig:teach-P-04-prototype-03}
\end{figure}

% END REVISION ADDITION teach-P-04-prototype-03
DynareMCP supplies ideas for investigative planning, source graphs and meaningful
differences between model versions. Its economic-model parser and numerical
solvers are not direct engines for SFC requirements. The reusable pattern is a
reviewable change supported by a dependency graph, with unresolved alternatives
preserved. A legal adapter would compare rules, definitions, scope and dates
rather than macroeconomic equations.

The proposed search use is equally specific. Historical DynareMCP investigation
records use UCB-style exploration to choose work; the currently inspected
MathDevMCP derivation controller explicitly describes a different bounded
evidence-ordering method, not MCTS. It would be inaccurate to label both as the
same proof-search engine \citep{localrepos}.

For LegalMath, a search node could contain a candidate interpretation and its
unresolved questions. Actions might retrieve a referenced FAQ, request a
counterexample, check a definition, or prepare a question for a compliance owner.
A search score may prioritize expected information gain or reduction in review
work. It cannot determine which interpretation is legally correct. A branch with
more permissive outcomes or more successful solver calls receives no special
authority. If several readings remain viable, the product retains the disagreement.

Start with a bounded work queue ordered by materiality and dependencies. MCTS or
UCB scheduling becomes a candidate only if the prototype reveals a substantial
search-allocation bottleneck. Its evaluation would compare useful questions
resolved per unit of review effort under the same budget, while correctness
continues to depend on source review and formal evidence. Transaction execution
uses approved rules, with no exploratory search in the decision path.
% END SOURCE UNIT P-04-prototype:03



% BEGIN SOURCE UNIT P-05-evidence:02
\subsection{A reproducible research collection}
\label{merge:P-05-evidence:02}


The papers are stored under \path{docs/papers/} using the requested filename
convention: the title, an underscore, the first author's surname and the year in
parentheses. Filesystem-unsafe title punctuation is normalized, while the manifest
retains the original title. Each
record includes its source URL, bibliographic identity, edition note, retrieval
date, byte count, PDF page count and SHA-256 digest. The library index distinguishes
distinct works from alternate editions and identifies discovery leads that were
not available as usable full text.

The technical notes identify which part of each paper supports a proposed
borrowing and what remains untested. They distinguish a semantics from its
implementation, a reported experiment from a reproduced result, and a publicly
available repository from a licensed and approved bank dependency. Full-text
extraction is retained for local search, but the source PDF remains the reference
where layout or equations matter. ResearchAssistant extraction is an aid to
inspection, not certification of the extracted text.

% BEGIN REVISION ADDITION teach-P-05-evidence-02
\begin{figure}[H]
\centering
\TeachingFlow{Named local document copy}{Exact passage and source hash}{Manuscript claim and support judgment}
\caption{An audit trail must show what was actually read for a citation.}\label{fig:teach-P-05-evidence-02}
\end{figure}

% END REVISION ADDITION teach-P-05-evidence-02
The SFC source snapshots and the earlier local-repository inspection records
remain under \path{.localresources/}, with their own manifest. The LaTeX and
bibliography are now canonical in \path{docs/monograph/}; the earlier
proposal entry point builds the same unified work. This arrangement permits another
reviewer to trace the proposal's substantive examples back to a document edition
without reconstructing the web searches. Redistributing the paper library is a
separate matter from using the local copies for this research; the source licenses
and access terms remain relevant.
% END SOURCE UNIT P-05-evidence:02

\section{Exact controller sequence}
\label{sec:controller-reference}

\begin{lstlisting}[caption={Main controller, with explicit terminal uncertainty.}]
start_run(source_packet, profile, required_policy):
    validate_source_identity_and_policy()
    commit(run = INITIALIZING, counters = 0)
    reserve_and_issue_initial_roles_with_global_budget()
    collect_until_complete_or_action_deadline()
    validate_outputs_and_freeze_initial_proposals()
    reconcile_inventory_candidates_and_required_roles()

    while has_unresolved_issue():
        if integrity_failure():
            close_as(FAILED_INTEGRITY); break
        if cancelled():
            close_as(CANCELLED); break
        if deadline_or_global_or_round_limit_reached():
            mark_remaining_issues_with_stop_reason()
            close_as(BLOCKED_UNRESOLVED); break
        if no_progress_limit_reached():
            mark_remaining_issues(NO_PROGRESS)
            close_as(BLOCKED_UNRESOLVED); break

        plan = choose_permitted_discriminating_actions()
        if plan.is_empty():
            mark_remaining_issues(HUMAN_REQUIRED)
            close_as(BLOCKED_UNRESOLVED); break
        commit_next_round_and_reserve_bounded_actions(plan)
        dispatch_only_reserved_actions()
        collect_results_until_each_action_is_terminal()
        validate_results_and_create_candidate_revisions()
        rerun_affected_checks_and_reconcile_all_open_issues()
        record_progress_signature()

    if not already_closed():
        apply_integrity_cancellation_and_deadline_guards()
    if not already_closed():
        check_required_roles_coverage_and_all_mandatory_checks()
        if any_requirement_incomplete():
            close_as(BLOCKED_UNRESOLVED)
        else:
            close_as(READY_FOR_REVIEW)
    write_complete_report_for_every_terminal_status()
    return report_reference
\end{lstlisting}

\section{Machine-readable blocked-result example}

\begin{lstlisting}[caption={A blocked result preserves alternatives and the next decision.}]
{
  "schema_version": "interpretation.v1",
  "run_status": "BLOCKED_UNRESOLVED",
  "source_ref": "23EC46-retained-version",
  "selected_slice": "paragraph-10",
  "candidate_ids": ["cashback-reimbursement", "cashback-fee-reduction"],
  "material_unresolved_issue_ids": ["classification-cashback"],
  "release_eligible": false,
  "processing_stop": "ROUND_LIMIT",
  "rounds_issued": 3,
  "probability_of_legal_correctness": null,
  "next_decision": "Establish the applicable fee-discount definition",
  "next_decision_owner_role": "LEGAL_INTERPRETATION_REVIEWER"
}
\end{lstlisting}

\section{Exact financial-rule serialization}
\label{sec:financial-serialization}

The complete financial-condition bundle is supplied in
\path{fixtures/spi-financial.bundle.json} beneath the specification directory.
It declares two money inputs, an unconditional scope for this isolated
subcondition, the expression, interpretation record and a verified span in
Annex 1 paragraph 3.1. The excerpt below is a complete expression node from
that bundle. Node identifiers make its trace reproducible.

\begin{minipage}{\textwidth}
\begin{lstlisting}[caption={The portfolio comparison in the supplied JSON bundle. Amounts are HK cents.}]
{"node_id":"portfolio.ge", "op":"compare", "cmp":"ge",
 "left":{"node_id":"portfolio.input", "op":"fact",
         "name":"portfolio"},
 "right":{"node_id":"portfolio.threshold", "op":"literal",
          "type":"money_hkd", "value":"4000000000"}}
\end{lstlisting}
\end{minipage}

The second comparison uses \texttt{net\_assets\_ex\_home} and
\texttt{"8000000000"}; an \texttt{any} node combines them. These inputs denote
amounts normalized under adopted wealth definitions, not any account balance
supplied under the same name. Release requires those definitions and applicability
to be reviewed. The narrow result name, \texttt{spi.financial}, identifies a
subcondition rather than a complete SPI assessment or trade authorization.


\section{Complete rule-language grammar}

An implementer can express the grammar below directly as tagged records. Every
\texttt{Expr} additionally carries a unique \texttt{node\_id}. A rule is
\texttt{(id, type, scope, body, interpretation\_id, source\_span\_ids)};
\texttt{scope} is Boolean and \texttt{body} has the declared result type.
\begin{lstlisting}[caption={Complete operator grammar for RuleIR 0.1; repeated arguments are ordered. The Java demonstration implements a stated subset.}]
Type = bool | integer | money_hkd | date
Expr = literal(Type, value)
     | fact(name) | rule(name)
     | all([Expr, ...]) | any([Expr, ...]) | not(Expr)
     | compare(eq | ge | gt, left, right)
     | add(left, right) | sub(left, right)
     | scale(arg, numerator, denominator)
     | if(condition, then_value, else_value)
     | default(base, [Exception, ...])
Exception = (exception_id, guard, value,
             interpretation_id, source_span_ids)
\end{lstlisting}


\section{Durable request accounting and worker recovery}

An action moves through \texttt{RESERVED}, \texttt{DISPATCHED}, and one of
\texttt{SUCCEEDED}, \texttt{FAILED}, \texttt{TIMED\_OUT}, or
\texttt{RESULT\_UNKNOWN}. The last state is necessary when the controller cannot
determine whether a remote service performed the action. A lost connection after
submission is different from a request rejected before dispatch. Both can lead
to missing results, but only the former may already have consumed remote work.

The reservation transaction records the action identifier, input hashes, issue
root, round, budget charge, adapter identity, timeout and idempotency key. An
outbox entry is written in the same transaction. A worker claims that entry with
a lease and a fencing token. The token increases when ownership changes, so an
old worker cannot commit a result after a newer owner has taken responsibility.
A valid late result can be retained as an observation without changing the closed
action's authoritative state.

The dispatch boundary requires care. If the outbox records only ``not yet sent''
and a worker crashes after sending but before updating the database, another
worker may send again. A provider idempotency key can make that retry refer to the
same operation, but not every provider supports this guarantee. The adapter must
state its capability. Without remote idempotency, a recovery policy can choose to
leave the result unknown rather than risk another paid request, or authorize a
new counted action. It cannot truthfully promise exactly once remote execution.

% BEGIN REVISION ADDITION teach-M06-17
% Teaching figure remains in the main chapter.
% END REVISION ADDITION teach-M06-17
The budget system reserves an upper bound on tokens or cost when those resources
are capped. Settlement records actual usage when known. A refund of unused
monetary reservation may be legitimate after a confirmed terminal response, but
it does not decrement the issued-action count. Otherwise a sequence of failed
requests could evade the maximum number of attempts. If actual usage exceeds a
provider's documented bound, the run records an accounting discrepancy and stops
new spending until the integration is investigated.

A compare-and-swap version protects issue updates. The worker returns evidence
against the issue revision it received. If the issue has changed, the controller
revalidates whether the evidence is still relevant before applying a resolution.
This prevents two simultaneous actions from overwriting each other's objections
or resolving an issue using a candidate that is no longer current.

Cancellation stops new reservations and attempts to cancel outstanding work.
Remote cancellation may be best effort. The report records outstanding or late
actions and keeps their budget charges. Cancellation cannot convert an incomplete
run into a clean result, and it cannot delete the source of a later adverse
finding. These details matter because regulatory evidence must survive ordinary
operational events, including a person closing a browser or restarting a worker.


% BEGIN SOURCE UNIT P-02-literature:01
\subsection{Scope and method of this review}
\label{merge:P-02-literature:01}


The review began with the seven requested links and expanded through their
references, forward-citation searches for Catala and Stipula, and searches on
executable law, legal reasoning, formalization, dates, testing and process
compliance. ResearchAssistant supplied discovery records and local PDF extraction.
The saved OpenAlex citation searches returned 52 entries for the selected Catala
record and 13 for the selected Stipula journal record. Those are retrieval counts
for particular indexed editions, not complete or comparable measures of influence.
Preprints, revised editions and new publications can be missing or duplicated.

The initial proposal retained 38 paper editions representing 37 distinct
works. The first unified library retained 50 editions representing 49 works,
including the later legal-search and uncertainty studies. The present
audit adds two legal-reliability studies, bringing the academic collection
to 52 editions representing 51 works. Regulatory documents, standards
and software records are retained alongside them.
The AAAI tax-reasoning paper and its extended preprint are retained separately;
all seven requested sources are present. Technical reading covered the relevant
methods, semantics, evaluation and supporting appendix material. The accompanying
notes identify the inspected sections, version differences, implementation checks
and limits of adoption \citep{libraryreview}. This is a focused technical review
supporting a product decision, rather than a claim to have enumerated every
citation. Unavailable secondary leads are recorded separately and do not support
implementation claims.

% BEGIN REVISION ADDITION teach-P-02-literature-01
% Teaching figure remains with the main discussion.
% END REVISION ADDITION teach-P-02-literature-01
The expanded search treats eight families separately: executable statutes;
defaults and defeasible priorities; obligations and violations; temporal process
compliance; authority and provenance; legal languages, compilers and testing;
neurosymbolic formalization and evaluation; and deployable rule engines. Important
foundations predate Catala, while contract languages have partially separate
citation communities. The retained forward-citation records have individual
dispositions in \path{docs/papers/coverage.md}; additional searches and backward
references supply the missing mechanisms. Selection favors primary sources that
define a needed mechanism or report an informative failure. Downloading a paper
or screening its abstract is distinguished from inspecting its technical method.
This is a bounded review of the prototype's design questions. Remaining retrieval
gaps are stated rather than counted as read papers.

Version control is particularly important here. The ARc paper was first posted
in 2025, but the inspected second version is from July 2026. The eFLINT paper's
inspected version is from August 2026. The Stipula preprint requested by the
sponsor has three authors; the later journal edition has a different author list.
Bibliographic records and local filenames make these distinctions explicit.
Published results are reported as their authors' findings; none of their
benchmarks or proof developments was independently rerun for this proposal.
% END SOURCE UNIT P-02-literature:01

\section{Release identity and reproducible builds}
\label{sec:release-build-record}

A build manifest records the bundle, generated source, runtime source, toolchain,
Java target, dependencies and JAR digest. A verification report then names the
build and exact JAR it tested. A release manifest names that build and verification
report, together with applicability and effective-time information. Approval
binds the final release manifest and bundle.

The direction matters. Embedding a manifest that contains the JAR's own digest
inside that JAR creates a circular identity problem. The current design keeps
these records external and acyclic. A reviewer approves the exact tested bytes
through the final manifest, rather than approving a mutable directory or a
human-readable version label.

Candidate compilation occurs before legal approval, because reviewers need a
concrete package and verification evidence to inspect. Compilation does not grant
deployment authority. Export into a controlled release route occurs only after
the required approvals and checks. This distinction lets engineers do useful
reversible work without repeatedly asking for permission, while preserving the
human decision at the point where it matters.

Reproducible builds strengthen the identity chain. Under the same declared
toolchain and inputs, build in separate directories with fixed archive-entry
timestamps and compare emitted source and binary hashes. A mismatch should be
investigated and recorded. Reproducibility does not prove semantic correctness,
but it reduces uncertainty about which bytes correspond to the reviewed source.

The local MVP uses synthetic identities and unsigned historical archives to
exercise workflow. A bank deployment needs enterprise identity, role separation,
key management and independently verified release signatures or equivalent
controls. A hash proves equality of bytes under its cryptographic assumption;
it does not prove who authorized those bytes. The manuscript's implementation
plan keeps that integration distinct from the local demonstration.
